\MMHostGuidanceHeading

\subsubsection{Ser Mestre de Jogo}
Como MJ, deves equilibrar duas prioridades:
\begin{itemize}
  \item O jogo deve ser divertido.
  \item O jogo deve ser justo.
\end{itemize}

Dá informações aos jogadores no momento certo. O contexto que a personagem conhece há muito tempo pode ser partilhado antecipadamente, idealmente pelo menos um dia antes da sessão. Entrega as informações recentes, sensíveis ou essenciais para o enredo mais tarde, quando se tornarem relevantes para o jogo.

\subsubsection{O desenrolar do jogo}
Procura dividir o jogo em fases:
\begin{enumerate}
  \item Apresentações, em que cada pessoa dá a conhecer a sua personagem.
  \item Um primeiro período de convívio, relações e pequenas pistas sem grandes consequências.
  \item Um acontecimento desencadeador que torne o mistério principal inequívoco.
  \item Investigação, consultas privadas com o MJ quando necessário e uma votação ou acusação final.
\end{enumerate}

\subsubsection{Escrever a história}
Começa pela resposta ao mistério principal: quem cometeu o ato, o que aconteceu ou qual é a verdadeira explicação sobrenatural ou oculta. Constrói a história de trás para a frente a partir dessa resposta. Liga cada personagem a pelo menos outra e dá a pelo menos metade delas informações que possam contribuir para resolver o mistério principal.

Para escrever a história, começa pelos factos individuais~\ref{subsec:atomic_facts}. São os factos vistos de fora. Lembra-te de que os jogadores podem mentir durante o jogo ou nunca descobrir alguns deles. Por isso, garante que, se um jogador tivesse acesso a todos estes factos, descobrir a verdade seria trivial.

Depois, escreve o contexto~\ref{subsec:background}. Também apresenta uma perspetiva externa e é, sobretudo, uma versão mais desenvolvida de todos os factos. Ao escrevê-lo, poderás querer voltar atrás e alterar os factos, à medida que ganhas uma nova perspetiva da história.

Enquanto escreves estas duas partes, cria uma versão completa das personagens, depois dos mistérios e, por fim, das pistas.

Só quando tiveres confiança no estado da história deves começar a escrevê-la da perspetiva de cada personagem. Nessas versões, podes embelezá-la, mostrar mais dos pensamentos e da personalidade de cada uma e acrescentar informações desnecessárias. Podes continuar a alterar toda a história, mas garante que os textos se mantêm coerentes entre si.

\subsubsection{Criar as personagens}
Lembra-te de que, à exceção das notas do MJ, o texto das personagens será mostrado aos jogadores. A maior parte deve ser escrita na perspetiva de cada personagem, enriquecida com as suas opiniões e personalidade. Deve haver sempre alguma informação extra que não seja necessária.

Tem em atenção a dificuldade das tarefas ao atribuir pontos.

\subsubsection{Equilíbrio do jogo}
Antes de pensar no equilíbrio, lembra-te de que o jogo deve ser divertido. Se os jogadores se estão a divertir a encontrar e resolver pistas enquanto discutem, está tudo bem.

Quanto à dificuldade, não tentes tornar o mistério excessivamente complexo. Não estás a escrever um livro nem a realizar uma série que os jogadores terão incontáveis horas para interpretar. Haverá muitas distrações.

Se alguém tivesse acesso a quase 100\% dos factos, a resposta ao mistério deveria ser óbvia. Com 75\% dos factos e meia hora para pensar, deveria ser bastante fácil resolver o mistério principal. Com 50\% dos factos e uma hora para pensar, deveria pelo menos ser possível chegar à resposta correta.

Tem em conta o tempo perdido com falsas pistas, o tempo que os jogadores passam a tentar descobrir os segredos uns dos outros e assim por diante.

À medida que as pistas são encontradas, a história construída pelos jogadores deve encaminhá-los naturalmente para a resolução do mistério principal.
